\documentclass[10pt,a4paper]{amsart}
\usepackage[margin=19mm]{geometry}
\usepackage{amsmath,amssymb,mathtools}
\usepackage[scr=rsfs,cal=euler]{mathalfa}
\usepackage{xfrac}
\usepackage[unicode,hidelinks,bookmarks=false]{hyperref}
\newcommand{\ie}{i.e.\ }
\newcommand{\cf}{cf.\ }
\newcommand{\resp}{respectively\ }
\newcommand{\Subsection}{\subsection}
\providecommand{\nobreakdash}{\nobreak\hbox{-}}
\setlength{\emergencystretch}{2em}
\multlinegap0pt
\usepackage{amssymb}
\usepackage{bm}
\usepackage[shortlabels]{enumitem}
\newtheorem{lemma}{Lemma}
\newcommand{\ASSharp}{{\mathcal{A}_{\sharp,\frac{7}{4}}^2}}
\newcommand{\al}{\alpha}
\newcommand{\bw}{\mathbf{W}}
\newcommand{\f}{\frac}
\newcommand{\p}{\partial}
\title{Low regularity well-posedness for two-dimensional hydroelastic waves}
\author{Lizhe Wan, Jiaqi Yang}
\date{Source: arXiv:2512.22040v1 (CC BY 4.0)}
\begin{document}
\maketitle

\noindent\emph{Source and licence.} Low regularity well-posedness for two-dimensional hydroelastic waves. Version arXiv:2512.22040v1 (CC BY 4.0), distributed by arXiv under the Creative Commons Attribution 4.0 International licence, http://creativecommons.org/licenses/by/4.0/, as recorded in the arXiv licence metadata for this exact version and shown on the abstract page https://arxiv.org/abs/2512.22040v1. Full licence notice: \texttt{licenses/arxiv-2512.22040-CC-BY-4.0.txt}.\par\smallskip

\noindent\textit{Editorial scope.} Counted unit: the article's lemma pexpansion (source0.tex lines 558--569). The article's complete original proof of it is delivered with it (source0.tex lines 571--655, one \texttt{proof} environment of 85 lines). No mathematics is written, repaired or re-derived: the statement and the proof are the article's own lines, in the article's own notation, and no retained line is substituted or re-flowed.  No further source-local context is needed: every cross-reference of every retained line resolves inside the two delivered spans.  Nothing was omitted from any retained span, so the package carries no omission record.  No retained line contains a citation, so the wrapper carries no bibliography and no bibliography entry.  No retained line contains a figure environment, an image-inclusion command or drawing code, so no image is delivered and the package declares no asset dependency.  Editorial formatting only, in addition: an AMS \texttt{amsart} preamble that restates the article's own declarations and macros used by the retained lines, namely the environments \texttt{lemma} on separate counters, together with the standard packages those lines need.  The retained environments print in this document as Lemma 1 in order of appearance, whereas the article prints the same items with the numbering of its own full document.  The complete native source of the article as distributed in the arXiv e-print for this exact version is bundled unchanged as \texttt{sources/191707/source0.tex}.
\begin{lemma}
The term $(1+\mathbf{W})p$ admits the following decomposition:
\begin{equation} \label{pexpansion}
\begin{split}
(1+\mathbf{W})p=&\p_{\al}(J^{-\f12}\p_{\al}(J^{-\f12}\p_{\al}(J^{-\f12}w_{\al})))\\
&-i\left[\p_{\al}\left(J^{-\f12}\p_{\al}\left(\tilde{c}J^{-\f12}w_{\al}\right)\right)
+\p_{\al}\left(\tilde{c}J^{-\f12}\p_{\al}\left(J^{-\f12}w_{\al}\right)\right)\right]\\
&-\f{11}4\p_{\al}\left(\tilde{c}^2J^{-\f12}w_{\al}\right)
-icw_{\al} + K.
\end{split}    
\end{equation}
\end{lemma}

\begin{proof}
Recall that $p$ is given by
\begin{align*}
p:=&\f{1}{J^{\f12}}\f{d}{d\al}\left[\f{1}{J^{\f12}}\f{d}{d\al}\left(\f{w_{\al\al}}{J^{\f12}(1+\mathbf{W})}-\left(\f{3\mathbf{W}_{\al}}{2J^{\f12}(1+\mathbf{W})^2}-\f{\mathbf{W}_{\al}}{2J^{\f32}}\right)w_{\al}\right)\right]\\
&-\f{1}{J^{\f12}}\f{d}{d\al}\left[\frac{(1+\bar{\mathbf{W}})w_{\al}}{2J^{\f32}}\f{d}{d\al}\left(\f{\mathbf{W}_{\al}}{J^{\f12}(1+\mathbf{W})}-\f{\bar{\mathbf{W}}_{\al}}{J^{\f12}(1+\bar{\mathbf{W}})}
\right)\right]\\
&-\frac{(1+\bar{\mathbf{W}})w_{\al}}{2J^{\f32}}\f{d}{d\al}\left[\f{1}{J^{\f12}}\f{d}{d\al}\left(\f{\mathbf{W}_{\al}}{J^{\f12}(1+\mathbf{W})}-\f{\bar{\mathbf{W}}_{\al}}{J^{\f12}(1+\bar{\mathbf{W}})}
\right)\right]\\
&+\f32\left(\f{w_{\al\al}}{J^{\f12}(1+\mathbf{W})}-\left(\f{3\mathbf{W}_{\al}}{2J^{\f12}(1+\mathbf{W})^2}-\f{\bar{\mathbf{W}}_{\al}}{2J^{\f32}}\right)w_{\al}\right)\left[\f{\mathbf{W}_{\al}}{J^{\f12}(1+\mathbf{W})}
-\f{\bar{\mathbf{W}}_{\al}}{J^{\f12}(1+\bar{\mathbf{W}})}\right]^2,
\end{align*}
which we decompose as
\begin{equation*}
   p :=A-B-C+D.
\end{equation*}
A direct computation shows that
\begin{align*}
\f{w_{\al\al}}{J^{\f12}(1+\mathbf{W})}-\left(\f{3\mathbf{W}_{\al}}{2J^{\f12}(1+\mathbf{W})^2}-\f{\bar{\mathbf{W}}_{\al}}{2J^{\f32}}\right)w_{\al}=\f{\p_{\al}(J^{-\f12}w_{\al})-i\tilde{c}w_{\al}}{1+\mathbf{W}}=:\f{A_0}{1+\mathbf{W}}.
\end{align*}
A straightforward calculation gives that
\begin{align*}
&J^{-\f12}\p_{\al}\left(\frac{A_0}{1+\mathbf{W}}\right)=\p_{\al}\left(\frac{J^{-\f12}A_0}{1+\mathbf{W}}\right)
-\frac{\p_{\al}J^{-\f12}A_0}{1+\mathbf{W}}\\
=&\frac{\p_{\al}\left(J^{-\f12}A_0\right)}{1+\mathbf{W}}-\f{\mathbf{W}_{\al}J^{-\f12}A_0}{(1+\mathbf{W})^2}
-\frac{\p_{\al}J^{-\f12}A_0}{1+\mathbf{W}}
=\f{\p_{\al}(J^{-\f12}A_0)-\f{i\tilde{c}A_0}{2}}{1+\mathbf{W}}=:\f{A_1}{1+\mathbf{W}}.
\end{align*}
Replace $A_0$ by $A_1$ in the above equations, we then get
\begin{align*}
A=&J^{-\f12}\p_{\al}\left(\frac{A_1}{1+\mathbf{W}}\right)=\f{\p_{\al}(J^{-\f12}A_1)-\f{i\tilde{c}A_1}{2}}{1+\mathbf{W}}\\
=&\f{1}{1+\mathbf{W}}\left(\p_{\al}(J^{-\f12}\p_{\al}(J^{-\f12}A_0))-\f12\p_{\al}\left((i\tilde{c})J^{-\f12}A_0\right)-\f{i\tilde{c}\p_{\al}(J^{-\f12}A_0)}{2}+\f{(i\tilde{c})^2A_0}{4}\right).
\end{align*}
Similarly, $B$ is given by
\begin{align*}
B=\f12J^{-\f12}\p_{\al}\left(\f{J^{-\f12}(i\tilde{c})_{\al}w_{\al}}{1+\mathbf{W}}\right)=\f12\f{\p_{\al}(J^{-\f12}J^{-\f12}(i\tilde{c})_{\al}w_{\al})-\f{i\tilde{c}J^{-\f12}(i\tilde{c})_{\al}w_{\al}}{2}}{1+\mathbf{W}}.
\end{align*}
The term $C$ is
\begin{align*}
C=\f12\f{1}{1+\mathbf{W}}J^{-\f12}\p_{\al}\left(J^{-\f12}(ic)_{\al}\right)w_{\al}.
\end{align*}
For $D$, we have
\begin{align*}
D=\f32\f{A_0}{1+\mathbf{W}}(i\tilde{c})^2.
\end{align*}
In the following, we will simplify terms in $(1+\bw)p$ according to the number of derivatives that may act on $w$.

{\bf{Fourth-order term of $w$ in $(1+\mathbf{W})p$}:}
\begin{equation*}
\p_{\al}(J^{-\f12}\p_{\al}(J^{-\f12}\p_{\al}(J^{-\f12}w_{\al}))).   
\end{equation*}

{\bf{Third-order terms of $w$ in $(1+\mathbf{W})p$:}}
\begin{align*}
&-\p_{\al}(J^{-\f12}\p_{\al}(J^{-\f12}(i\tilde{c})w_{\al}))
-\f12\p_{\al}\left((i\tilde{c})J^{-\f12}\p_{\al}(J^{-\f12}w_{\al})\right)
-\f{i\tilde{c}\p_{\al}(J^{-\f12}\p_{\al}(J^{-\f12}w_{\al}))}{2}\\
=&-\p_{\al}(J^{-\f12}\p_{\al}(J^{-\f12}(i\tilde{c})w_{\al}))
-\p_{\al}\left((i\tilde{c})J^{-\f12}\p_{\al}(J^{-\f12}w_{\al})\right)
+\f{(i\tilde{c})_{\al}J^{-\f12}\p_{\al}(J^{-\f12}w_{\al})}{2}.
\end{align*}

{\bf{Second-order terms of $w$ in $(1+\mathbf{W})p$:}}
\begin{align*}
&\f12\p_{\al}\left((i\tilde{c})^2J^{-\f12}w_{\al}\right)+\f{i\tilde{c}\p_{\al}(J^{-\f12}(i\tilde{c})w_{\al})}{2}+\f{(i\tilde{c})^2\p_{\al}(J^{-\f12}w_{\al})}{4}\\
&+\f{3(i\tilde{c})^2\p_{\al}(J^{-\f12}w_{\al})}{2}+\f{(i\tilde{c})_{\al}J^{-\f12}\p_{\al}(J^{-\f12}w_{\al})}{2}-\f12\p_{\al}(J^{-\f12}J^{-\f12}(i\tilde{c})_{\al}w_{\al})\\
=&\f{11}4\p_{\al}\left((i\tilde{c})^2J^{-\f12}w_{\al}\right)-4(i\tilde{c})_{\al}(J^{-\f12}(i\tilde{c})w_{\al})
-\f12J^{-\f12}\p_{\al}(J^{-\f12}(i\tilde{c})_{\al})w_{\al}.
\end{align*}


{\bf{First-order terms of $w$ in $(1+\mathbf{W})p$:}}
\begin{align*}
&-4(i\tilde{c})_{\al}(J^{-\f12}(i\tilde{c})w_{\al})+\f14(i\tilde{c})_{\al}(J^{-\f12}(i\tilde{c})w_{\al}) -\f12J^{-\f12}\p_{\al}(J^{-\f12}(i\tilde{c})_{\al})w_{\al}
\\
&-\f12J^{-\f12}\p_{\al}(J^{-\f12}(i\tilde{c})_{\al})w_{\al}
-\f14(i\tilde{c})^3w_{\al}-\f32(i\tilde{c})^3w_{\al}\\
=&-icw_{\al}+\f{15}{4}\tilde{c}_{\al}\tilde{c}J^{-\f12}w_{\al}-\f{5i}{4}\tilde{c}^3w_{\al}.
\end{align*}
Note that 
\begin{equation*}
 \left\|\f{15}{4}\tilde{c}_{\al}\tilde{c}J^{-\f12}w_{\al}-\f{5i}{4}\tilde{c}^3w_{\al} \right\|_{L^2} \lesssim \ASSharp \| w\|_{H^\f32} ,
\end{equation*}
so that this term can be put into the perturbative source term $K$.
Collecting the contributions at each derivative order in $(1+\bw)p$, we obtain the expansion in \eqref{pexpansion}. 
\end{proof}

\end{document}
